\documentclass[12pt]{article}

% ============================================================
%  BEER manuscript — full LaTeX source
%  Biochemical Estimator & Explorer of Residues, v2
%  Author: Saumyak Mukherjee
% ============================================================

\usepackage[utf8]{inputenc}
\usepackage[T1]{fontenc}
\usepackage{lmodern}
\usepackage[margin=2.5cm]{geometry}
\usepackage{amsmath,amssymb}
\usepackage{graphicx}
\usepackage{booktabs}
\usepackage{longtable}
\usepackage{array}
\usepackage{xcolor}
\usepackage{hyperref}
\usepackage{natbib}
\usepackage{setspace}
\usepackage{microtype}
\usepackage{caption}
\usepackage{enumitem}
\usepackage{authblk}

\hypersetup{
  colorlinks=true,
  linkcolor=blue!70!black,
  citecolor=green!50!black,
  urlcolor=blue!70!black,
}

\onehalfspacing

% ============================================================
\title{\textbf{BEER: Biochemical Estimator \& Explorer of Residues}\\[6pt]
\large A Comprehensive, Cross-Platform Desktop Suite for Protein\\
Sequence and Structure Analysis}

\author[1]{Saumyak Mukherjee}
\affil[1]{Department of Theoretical Biophysics,
          Max Planck Institute of Biophysics,\\
          Max-von-Laue-Str.~3, 60438 Frankfurt am Main, Germany}
\date{}

% ============================================================
\begin{document}
\maketitle

\begin{abstract}
Rapid, in-depth biochemical characterisation of protein sequences is a daily
requirement in structural biology, proteomics, and the burgeoning field of
intrinsically disordered proteins (IDPs).
Existing web servers are powerful but are session-bound, require Internet
connectivity, and typically address only one or two analytical domains at
a time.
Here we present \textbf{BEER} (Biochemical Estimator \& Explorer of Residues),
a cross-platform, open-source desktop application that integrates twenty-three
distinct analysis categories and thirty-one interactive, publication-quality
visualisations into a single, self-contained Python script.
Beyond classical physicochemical properties (molecular weight, isoelectric
point, extinction coefficient, GRAVY, instability index), BEER provides
advanced modules covering: $\beta$-aggregation propensity (ZYGGREGATOR/PASTA),
intrinsic solubility (CamSol), post-translational modification (PTM)
scanning, signal peptide and GPI-anchor prediction (von~Heijne model),
amphipathic helix detection (Eisenberg hydrophobic moment), sequence charge
decoration (SCD, Sawle \& Ghosh 2015), RNA-binding propensity, tandem repeat
identification, eukaryotic linear motif (ELM) database queries, DisProt
disorder annotations, PhaSepDB phase-separation look-up, and
Ramachandran/contact-network analysis of experimental or AlphaFold structures.
Sequences or structures may be entered manually, imported from FASTA/PDB
files, or fetched from UniProt and RCSB PDB by accession number, with
automatic post-fetch analysis.
A multichain comparison module, an NCBI BLAST interface, Pfam domain
architecture display, and an integrated 3D structure viewer complete the
feature set.
All computations are performed locally using standard scientific Python
libraries; no registration or API key is required.
BEER runs on Windows, macOS, and Linux and is distributed as a single
Python file (\textasciitilde{}10,000 lines) under the GNU GPLv2 licence.
\end{abstract}

\noindent\textbf{Keywords:} protein sequence analysis; intrinsically disordered
proteins; phase separation; post-translational modifications; bioinformatics;
graphical user interface; AlphaFold; BLAST

% ============================================================
\section{Introduction}
\label{sec:intro}

Proteins are the workhorses of biology, and understanding their physicochemical
properties is fundamental to predicting their function, stability, and
interaction partners.
Over the past two decades a rich ecosystem of web-based tools has emerged to
address this need: ProtParam~\citep{gasteiger2005} computes basic sequence
properties; IUPred~\citep{dosztanyi2005} and PONDR~\citep{romero2001} predict
intrinsic disorder; Zyggregator~\citep{tartaglia2008} assesses aggregation
propensity; CamSol~\citep{sormanni2015} quantifies intrinsic solubility;
SignalP~\citep{almagro2019} identifies signal peptides; and the ELM
database~\citep{kumar2020} catalogues eukaryotic linear motifs.
Each of these tools is authoritative in its domain, but the practitioner who
needs a holistic picture of a newly discovered protein must visit five or more
separate servers, copy-paste sequences repeatedly, and reconcile results
presented in incompatible formats.

Furthermore, web servers are inherently session-limited: results are not
permanently stored, analyses cannot be run offline, and large-scale or
repeated queries are constrained by server-side rate limits.
For protein biochemists working with proprietary or pre-publication sequences,
the requirement to upload data to external servers raises legitimate data
security concerns.

Desktop applications address these shortcomings but have historically
lagged behind web tools in analytical depth.
Here we describe BEER (Biochemical Estimator \& Explorer of Residues), a
comprehensive desktop application that brings together, in a single Python
script, the analytical capabilities previously scattered across dozens of
servers and scripts.
BEER is built on standard, freely available scientific Python libraries
(BioPython, NumPy, Matplotlib, PyQt5) and therefore requires no proprietary
software, no registration, and no network connectivity for its core analyses.
Network-dependent features (BLAST, UniProt/PDB fetch, ELM, DisProt,
PhaSepDB, AlphaFold, Pfam) are clearly separated and fail gracefully in
offline environments.

The initial release of BEER (v1, \citealt{mukherjee2025}) established the
core physicochemical analysis and visualisation framework.
The present work describes BEER v2, which adds eight new analytical modules
(aggregation, PTM scanning, signal peptide prediction, amphipathic helix
detection, SCD, RNA-binding propensity, tandem repeat detection, and the
new graphs sub-system), five network-retrieval workers (ELM, DisProt,
PhaSepDB, AlphaFold, Pfam), full structural analysis (Ramachandran plots,
Ca distance maps, contact networks, 3D viewer), an auto-analysis workflow,
a multichain comparison module, and extensive usability improvements
including a sidebar navigation, a dark-mode toggle, and per-graph export at
up to 200~dpi.

% ============================================================
\section{Implementation}
\label{sec:impl}

\subsection{Language, dependencies, and distribution}

BEER is written in Python~3.11--3.12 and distributed as a single executable
file (\texttt{beer.py}).
All analytical sub-modules have been inlined into this file to eliminate
installation complexity; no package structure is required beyond the
following pip-installable dependencies:

\begin{center}
\begin{tabular}{ll}
\toprule
\textbf{Package} & \textbf{Purpose} \\
\midrule
BioPython $\geq$1.80 & Sequence I/O, ProteinAnalysis, BLAST, PDB parsing \\
NumPy $\geq$1.24 & Numerical arrays \\
Matplotlib $\geq$3.7 & Publication-quality plots \\
PyQt5 $\geq$5.15 & Cross-platform GUI \\
mplcursors & Interactive hover cursors on plots \\
PyQtWebEngine (optional) & Embedded 3-D structure viewer (NGL) \\
\bottomrule
\end{tabular}
\end{center}

The optional PyQtWebEngine dependency activates the interactive 3-D structure
viewer; all other features are fully available without it.

\subsection{Architecture}

BEER follows a three-tier architecture: (1) a set of pure Python computation
functions that operate entirely on strings and NumPy arrays and carry no GUI
dependencies; (2) asynchronous Qt worker threads (\texttt{QThread} subclasses)
that wrap network-IO operations and emit PyQt signals on completion; and (3) a
\texttt{MainWindow} class that assembles the GUI, connects signals, and
delegates analysis to the compute layer.

This separation allows the computation layer to be imported and called from
automated test scripts without instantiating a display, as verified by the
unit tests in \texttt{tests/}.

\subsection{User interface}

The main window is organised around a left-side icon navigation bar containing
eight sections (Figure~\ref{fig:screenshot}):

\begin{enumerate}[noitemsep]
  \item \textbf{Analysis} — sequence input, chain selector, and thirteen
        scrollable report sections.
  \item \textbf{Graphs} — a tree-style browser giving access to all 31 plots
        grouped by category.
  \item \textbf{Structure} — 3-D viewer, Ramachandran plot, Ca distance map,
        and contact network.
  \item \textbf{BLAST} — NCBI blastp against nr, SwissProt, PDB, or
        RefSeq\_protein.
  \item \textbf{Compare} — side-by-side physicochemical comparison of two
        sequences.
  \item \textbf{Multichain Analysis} — batch summary table with CSV/JSON
        export.
  \item \textbf{Settings} — all configurable parameters.
  \item \textbf{Help} — built-in method and graph reference.
\end{enumerate}

Analysis of a newly loaded sequence (from file, manual entry, or API fetch)
is triggered automatically, eliminating the need for an explicit ``Analyse''
button click in the common workflow.

% ============================================================
\section{Analysis Modules}
\label{sec:modules}

\subsection{Composition and basic physicochemical properties}

Amino acid composition (count and percentage frequency for all 20 standard
residues) is computed directly from the sequence.
Molecular weight, isoelectric point (pI), GRAVY score, instability index,
aromaticity, and extinction coefficient at 280~nm are delegated to BioPython's
\texttt{ProteinAnalysis} class~\citep{cock2009}.
Net charge as a function of pH is computed using the Henderson--Hasselbalch
equation with user-overridable pKa values:

\begin{equation}
Q(\mathrm{pH}) =
  \sum_{i \in \text{basic}} \frac{1}{1 + 10^{\mathrm{pH} - \mathrm{p}K_a^{(i)}}}
  - \sum_{j \in \text{acidic}} \frac{1}{1 + 10^{\mathrm{p}K_a^{(j)} - \mathrm{pH}}}
\end{equation}

Default pKa values are: N-terminus 9.69, C-terminus 2.34, Asp 3.90, Glu 4.07,
Cys 8.18, Tyr 10.46, His 6.04, Lys 10.54, Arg 12.48.

\subsection{Hydrophobicity}

Per-residue hydrophobicity profiles are computed with the Kyte--Doolittle
scale~\citep{kyte1982} using a sliding window whose width is configurable
(default 9 residues) from the Settings panel.
The GRAVY score (Grand Average of hYdropathicity) is the arithmetic mean of
per-residue KD values.
Residues are additionally classified as hydrophobic (KD\,>\,0),
hydrophilic (KD\,<\,0), or neutral (KD\,=\,0).

\subsection{Charge analysis}

Four charge descriptors are computed for every sequence:

\begin{itemize}[noitemsep]
  \item \textbf{FCR} (fraction of charged residues): $(n_+ + n_-) / N$.
  \item \textbf{NCPR} (net charge per residue): $(n_+ - n_-) / N$.
  \item \textbf{$\kappa$} (charge patterning): the Das--Pappu parameter that
        measures the degree to which positive and negative charges are
        segregated along the chain; $\kappa = 0$ indicates a well-mixed
        sequence and $\kappa = 1$ indicates maximum segregation~\citep{das2013}.
  \item \textbf{$\Omega$} (charge-aromatic patterning): extension of $\kappa$
        to a combined set of charged and aromatic residues~\citep{das2015}.
\end{itemize}

The sliding-window local charge profile (NCPR vs position) is displayed as
a filled line plot with distinct positive and negative fill regions.

\subsection{Aromatic and $\pi$-interactions}

Aromatic fraction ($f_{\text{arom}} = (n_Y + n_F + n_W) / N$), cation--$\pi$
pairs (K/R paired with F/W/Y within $\pm$4 residues), and $\pi$--$\pi$ pairs
(F/W/Y paired with F/W/Y within $\pm$4 residues) are computed.
The cation--$\pi$ heat map provides a 2-D visualisation with colour intensity
proportional to $1/d$ where $d$ is the sequence distance between the pair.

\subsection{Low-complexity and Shannon entropy}

Sequence complexity is quantified by the Shannon entropy:

\begin{equation}
H = -\sum_{i=1}^{20} f_i \log_2 f_i
\end{equation}

where $f_i$ is the frequency of amino acid $i$.
The normalised entropy $H/\log_2(20)$ ranges from 0 (single residue type)
to 1 (perfectly uniform composition).
Low-complexity windows (w\,=\,12, $H < 2.0$\,bits) are identified and
their fraction reported.
A prion-like score is computed as the fraction of N, Q, S, G, and Y
residues~\citep{lancaster2014}.

\subsection{Disorder and secondary structure propensity}

A disorder propensity profile is computed using an IUPred-inspired scoring
scheme: disorder-promoting residues (A, E, G, K, P, Q, R, S) contribute
positively and order-promoting residues (C, F, H, I, L, M, V, W, Y)
contribute negatively, smoothed over a sliding window.
The aliphatic index (Ikai 1980) and the Das--Pappu $\Omega$ parameter are
also reported.

Secondary structure propensities are derived from the Chou--Fasman
tables~\citep{chou1974}: per-residue $P_\alpha$ (helix) and $P_\beta$ (sheet)
propensities are averaged over the sliding window and displayed as two-panel
filled plots.

\subsection{Transmembrane helix prediction}

Transmembrane (TM) helices are predicted using a re-implementation of the
Kyte--Doolittle sliding-window method~\citep{kyte1982} with a window of
19 residues and a threshold of 1.6.
Topology is assigned using the von~Heijne inside-positive rule: the
flanking side with more K/R residues is designated as cytoplasmic.
Contiguous segments of 17--25 residues are retained as TM helices; overlength
segments are trimmed to the single highest-scoring 19-residue window.

\subsection{Phase-separation propensity and LARKS}

A composite liquid--liquid phase-separation (LLPS) propensity score
($0$--$1$) is computed as a weighted linear combination of seven features:
aromatic fraction ($w = 0.25$), prion-like score ($w = 0.20$), disorder
fraction ($w = 0.20$), FCR ($w = 0.10$), $\Omega$ ($w = 0.10$),
normalised LARKS density ($w = 0.15$), and an $|\mathrm{NCPR}|$ penalty
($w = -0.10$).
LARKS (low-complexity amyloid-like reversible kinked segments)~\citep{hughes2018}
are identified as 7-residue windows with $\geq$1 aromatic residue, $\geq$50\%
low-complexity residues, and Shannon entropy\,$< 1.8$\,bits.

\subsection{Linear motifs}

A built-in regex scanner searches for 15 short linear motifs (SLiMs):
nuclear localisation signal (NLS), nuclear export signal (NES), PxxP
(SH3 ligand), 14-3-3 binding, RGG box, FG repeat, KFERQ (chaperone-mediated
autophagy), KDEL (ER retention), PKA site (RxxS/T), SxIP (EB1 ligand), WW
domain ligand, caspase-3 cleavage site, N-glycosylation sequon, SUMOylation
consensus (ψKxE), and CK2 site ([ST]xx[DE]).

For UniProt accessions an asynchronous \texttt{ELMWorker} thread queries
the Eukaryotic Linear Motif (ELM) database REST API~\citep{kumar2020} to
retrieve experimentally validated motif instances, complementing the
regex-based predictions.

\subsection{Aggregation and solubility}

The aggregation module provides three complementary quantitative metrics:

\begin{enumerate}[noitemsep]
  \item \textbf{ZYGGREGATOR profile}: a 6-residue sliding-window average of
        the Tartaglia--Vendruscolo per-residue $\beta$-aggregation propensity
        scale~\citep{tartaglia2008}.
        Hotspots are defined as contiguous stretches of $\geq$4 residues with
        mean propensity $\geq 1.0$.
  \item \textbf{PASTA energy}: for each hotspot the mean diagonal PASTA pairwise
        $\beta$-strand interaction energy~\citep{trovato2007} is reported
        (more negative = stronger amyloidogenic tendency).
  \item \textbf{CamSol intrinsic score}: per-residue intrinsic solubility from
        the CamSol scale~\citep{sormanni2015}.
        The mean score, fraction of insoluble residues (CamSol\,$< -0.2$),
        and fraction of soluble residues (CamSol\,$> 0.2$) are reported, along
        with a 7-residue moving-average smoothed profile.
\end{enumerate}

\subsection{Post-translational modification prediction}

The PTM scanner (\texttt{scan\_ptm\_sites}) identifies ten types of
putative modification sites using published consensus sequence rules:

\begin{itemize}[noitemsep]
  \item N-linked glycosylation: N[^P][ST] sequon (high confidence).
  \item O-linked glycosylation: $\geq$3 S/T in any 5-aa window (mucin-type;
        medium confidence).
  \item Phosphoserine/Thr (CK2): [ST]xx[DE] (medium confidence).
  \item Phosphoserine/Thr (PKA): R[^P][^P][ST] (medium confidence).
  \item Phosphotyrosine (EGFR-like): Y flanked by D/E within $\pm$3
        (low confidence).
  \item Ubiquitination (ψKXE): [LVIMF]K.[DE] (medium confidence).
  \item SUMOylation (ψKxE): [VILMF]K.E (medium confidence).
  \item N-terminal acetylation (NatA): Met removal rule or N-terminal S/A/T/G/C.
  \item Lysine acetylation: KxxK or GKxx motifs (low confidence).
  \item Arginine methylation: RGG, RG, and GR motifs (medium confidence).
\end{itemize}

Each predicted site is reported with its position (1-based), flanking sequence
context ($\pm$3 residues), a textual description, and a confidence tier
(high/medium/low).

\subsection{Signal peptide and GPI-anchor prediction}

Signal peptide prediction follows the von~Heijne three-region (n-h-c) model:

\begin{enumerate}[noitemsep]
  \item An \emph{n-region} (1--5 positively charged residues) is identified at
        the N-terminus.
  \item An \emph{h-region} of 10--15 residues is located immediately downstream
        and must achieve a mean KD score exceeding 1.4.
  \item A \emph{c-region} (5--7 residues) containing the cleavage site is
        required to have small neutral residues at positions $-$3 and $-$1
        relative to the cleavage site (AXA rule, von~Heijne 1983).
\end{enumerate}

The full signal peptide is reported with its predicted cleavage site,
subcellular localisation, and a confidence score based on h-region KD average.
GPI-anchor prediction identifies a C-terminal hydrophobic region ($\geq$8
residues, mean KD\,$\geq 1.2$) preceded by a small omega-site amino acid
(A/S/T/D/N/G/C), consistent with GPI-anchor addition rules.

\subsection{Amphipathic helix detection}

Amphipathic helices are detected using the Eisenberg hydrophobic
moment~\citep{eisenberg1984}:

\begin{equation}
\mu_H = \frac{1}{n}\sqrt{\left(\sum_{i=0}^{n-1} H_i \sin(i\delta)\right)^2
                        + \left(\sum_{i=0}^{n-1} H_i \cos(i\delta)\right)^2}
\end{equation}

where $H_i$ is the Eisenberg consensus hydrophobicity of residue $i$,
$\delta = 100°$ for $\alpha$-helices (or $160°$ for $\beta$-strands),
and $n$ is the window length (default 11 residues).
A helix is classified as amphipathic when $\mu_H \geq 0.38$
(the threshold of Eisenberg et al.\ 1984) and the mean hydrophobicity
$\langle H \rangle \geq 0.1$.
The membrane-active, surface-seeking, and globular helix classification
follows the original Eisenberg--Weiss diagram.

\subsection{Sequence charge decoration (SCD)}

The Sawle--Ghosh SCD parameter~\citep{sawle2015} is computed as:

\begin{equation}
\mathrm{SCD} = \frac{1}{N} \sum_{i < j} \sigma_i \sigma_j |i - j|^{1/2}
\end{equation}

where $\sigma_i = +1$ for K/R, $\sigma_i = -1$ for D/E, and $\sigma_i = 0$
otherwise.
Positive SCD indicates clustering of like charges; negative SCD indicates
polyampholyte-like alternating charge patterning, relevant to IDP phase
behaviour.
A sliding-window SCD profile (window\,=\,20) is also computed.
Additional block-length statistics (mean block length of K/R and D/E) and a
charge segregation score are reported.

\subsection{RNA-binding propensity}

RNA-binding propensity is quantified using a two-component score:

\begin{equation}
S_{\mathrm{RBP}} = 0.6 \cdot \min\!\left(\frac{\langle p \rangle + 0.3}{1.0}, 1\right)
               + 0.4 \cdot \min\!\left(\frac{f_{\mathrm{RBP}}}{0.25}, 1\right)
\end{equation}

where $\langle p \rangle$ is the mean per-residue propensity from the
Jeong et al.\ 2012 scale and $f_{\mathrm{RBP}}$ is the fraction of
high-propensity RNA-binding residues (K, R, Y, F, W, G, H).
A panel of eight RNA-binding motif patterns (RGG box, RG repeats, KH domain
GXXG loop, SR repeats, YGG/GGY, RRM RNP1, C$_2$H$_2$ zinc finger, and
DEAD/DEAH box) is scanned by regular expression, and each match is reported
with its position and sequence.

\subsection{Tandem repeat identification}

Exact tandem repeats are detected by a greedy k-mer extension algorithm
operating over unit lengths of 2--8 residues and requiring at least 2
consecutive copies.
The output reports start/end positions (1-based), repeat unit sequence, unit
length, copy number, and total covered length.
Overlapping candidates at the same position are resolved by retaining the
longest repeat.

\subsection{Coiled-coil prediction}

A coiled-coil propensity profile is computed using a 28-residue sliding window
and the heptad repeat periodicity score.
The HEPT score measures how well the hydrophobic residues at heptad positions
$a$ and $d$ (using a 3.5-residue periodicity) match the expected pattern,
with coiled-coil regions predicted where the score exceeds 0.50.

\subsection{Structural analysis from PDB/AlphaFold coordinates}

When atomic coordinates are available (from direct PDB file import, RCSB PDB
fetch, or AlphaFold EBI download), BEER provides four structural analyses:

\begin{enumerate}[noitemsep]
  \item \textbf{Ramachandran plot}: $\phi$/$\psi$ dihedral angles extracted
        via BioPython's \texttt{PPBuilder}; residues are colour-coded by
        secondary structure class (helix, sheet, coil) estimated from the
        angular region occupied.
  \item \textbf{C$\alpha$ distance map}: pairwise Euclidean distances between
        C$\alpha$ atoms rendered as a heatmap with an 8\,\AA\ contact contour
        overlaid.
  \item \textbf{Contact network}: graph-based visualisation where nodes
        represent residues and edges connect residues within 8\,\AA\; node
        colour encodes residue type and degree.
  \item \textbf{pLDDT / B-factor profile}: AlphaFold per-residue confidence
        scores (stored in the B-factor column) plotted as a filled profile;
        threshold lines at 70 and 90 delineate low/moderate/high confidence
        regions.
\end{enumerate}

An embedded NGL 3-D viewer (requires PyQtWebEngine) supports rotation,
zooming, and surface/ribbon/ball-and-stick display modes.
All three structure sources (manual PDB import, RCSB fetch, AlphaFold
download) share the same visualisation pipeline.

\subsection{External database integrations}

Four asynchronous worker threads provide optional network-dependent enrichment:

\begin{description}[noitemsep]
  \item[ELMWorker] Queries the ELM REST API
        (\texttt{https://elm.eu.org/instances.json?q=\{accession\}}) to
        retrieve experimentally validated eukaryotic linear motif instances.
  \item[DisPRotWorker] Queries the DisProt REST API
        (\texttt{https://disprot.org/api/\{accession\}}) to retrieve
        experimentally characterised disordered regions.
  \item[PhaSepDBWorker] Queries PhaSepDB
        (\texttt{https://phasepdb.org/api/protein/\{accession\}}) to
        retrieve experimental phase-separation data (category, evidence type,
        organism, references).
  \item[AlphaFoldWorker] Downloads the AlphaFold2 predicted structure from
        the EBI server and the associated pLDDT JSON file; the retrieved
        structure is loaded directly into the Structure section for immediate
        visualisation.
\end{description}

The BLAST module submits sequences to NCBI via \texttt{Bio.Blast.NCBIWWW}
(blastp, configurable E-value cutoff) and parses the XML results into a
sortable hit table; selecting any hit re-analyses its sequence directly within
BEER.
Pfam domain annotations are retrieved from the EMBL-EBI InterPro REST API
and rendered as coloured domain blocks in the Domain Architecture graph.

% ============================================================
\section{Visualisations}
\label{sec:viz}

BEER produces 31 publication-quality interactive plots rendered at 120\,dpi
on-screen and 200\,dpi on export.
All plots embed \texttt{mplcursors} hover annotations.
The complete catalogue is summarised in Table~\ref{tab:graphs}.

\begin{longtable}{p{4.5cm}p{9.5cm}}
\caption{BEER graph catalogue (31 plots).}
\label{tab:graphs}\\
\toprule
\textbf{Graph} & \textbf{Description} \\
\midrule
\endfirsthead
\toprule
\textbf{Graph} & \textbf{Description} \\
\midrule
\endhead
\midrule
\multicolumn{2}{r}{\itshape continued\ldots}\\
\endfoot
\bottomrule
\endlastfoot
%
\multicolumn{2}{l}{\textit{Composition}}\\
AA Composition (Bar)         & Residue counts with frequency annotations, sorted by frequency, hydrophobicity, or A--Z \\
AA Composition (Pie)         & Residue frequencies; zero-count residues hidden \\
\midrule
\multicolumn{2}{l}{\textit{Profiles}}\\
Hydrophobicity Profile       & Sliding-window KD average; configurable window \\
Local Charge Profile         & Sliding-window NCPR; positive/negative fill \\
Local Complexity             & Sliding-window Shannon entropy; low-complexity threshold line \\
Disorder Profile             & IUPred-inspired per-residue score; orange fill above 0.5 \\
Coiled-Coil Profile          & Heptad-periodicity score; fill above 0.50 marks predicted coils \\
Linear Sequence Map          & Four-track view: hydrophobicity, NCPR, disorder, helix $P_\alpha$ \\
Secondary Structure          & Two-panel: helix $P_\alpha$ (top) and sheet $P_\beta$ (bottom); fill above 1.0 \\
\midrule
\multicolumn{2}{l}{\textit{Charge \& $\pi$-interactions}}\\
Net Charge vs pH             & Henderson--Hasselbalch curve pH 0--14; pI annotated \\
Isoelectric Focus            & Enhanced curve with pH 7.4 annotation \\
Charge Decoration            & Das--Pappu FCR vs $|$NCPR$|$ phase diagram with IDP boundary \\
Cation--$\pi$ Map            & K/R $\leftrightarrow$ F/W/Y proximity heat map ($\pm$8 residues) \\
\midrule
\multicolumn{2}{l}{\textit{Structure \& folding}}\\
Bead Model (Hydrophobicity)  & Per-residue KD scatter; 30+ colourmap choices \\
Bead Model (Charge)          & K/R blue, D/E red, H cyan, neutral grey \\
Sticker Map                  & Aromatic amber, basic blue, acidic pink, spacer grey \\
Helical Wheel                & 18-residue Cartesian projection; connecting lines; contrast labels \\
TM Topology                  & Snake-plot of predicted TM helices with inside-positive topology \\
\midrule
\multicolumn{2}{l}{\textit{Structural (requires coordinates)}}\\
pLDDT / B-factor Profile     & Per-residue AlphaFold confidence; threshold lines at 70 and 90 \\
C$\alpha$ Distance Map       & Pairwise distance heatmap; 8\,\AA\ contact contour \\
Ramachandran Plot            & $\phi$/$\psi$ scatter colour-coded by secondary structure \\
Contact Network              & Residue contact graph; node colour by type and degree \\
Domain Architecture          & Multi-track: Pfam domains, disorder, low-complexity, TM helices \\
\midrule
\multicolumn{2}{l}{\textit{IDP / phase separation}}\\
Uversky Phase Plot           & Mean $|$net charge$|$ vs mean hydrophobicity; IDP/ordered boundary \\
Saturation Mutagenesis       & 20$\times n$ $|\Delta\mathrm{GRAVY}| + |\Delta\mathrm{NCPR}|$ heatmap ($\leq$500\,aa) \\
\midrule
\multicolumn{2}{l}{\textit{Advanced analysis graphs}}\\
Aggregation Profile          & ZYGGREGATOR sliding-window profile with hotspot shading \\
CamSol Solubility Profile    & Smoothed per-residue CamSol intrinsic score \\
SCD Profile                  & Sliding-window SCD ($w=20$) \\
RBP Profile                  & Sliding-window RNA-binding propensity \\
Amphipathic Helix Moment     & Per-residue Eisenberg $\mu_H$ profile with amphipathic regions shaded \\
\end{longtable}

Graphs can be exported individually (PNG/SVG/PDF at 200\,dpi) or in batch
(Save All Graphs to a directory).
Transparent background export is enabled by default.

% ============================================================
\section{Settings and customisation}

BEER exposes eighteen user-configurable parameters through the Settings section
(Table~\ref{tab:settings}).
All settings take effect immediately upon clicking \emph{Apply Settings};
factory defaults are restored by \emph{Reset to Defaults}.

\begin{longtable}{p{4cm}p{2.5cm}p{7.5cm}}
\caption{Configurable BEER settings.}
\label{tab:settings}\\
\toprule
\textbf{Setting} & \textbf{Default} & \textbf{Description} \\
\midrule
\endfirsthead
\toprule
\textbf{Setting} & \textbf{Default} & \textbf{Description} \\
\midrule
\endhead
\bottomrule
\endlastfoot
Default pH                  & 7.0    & pH for net-charge calculations \\
Sliding window size         & 9      & Window for hydrophobicity, charge, entropy profiles \\
Override pKa                & ---    & Per-residue pKa override (comma-separated) \\
Reducing conditions         & off    & Free Cys for extinction coefficient \\
Label font size             & 14\,pt & Axis label font size \\
Tick font size              & 12\,pt & Tick label font size \\
Default graph format        & PNG    & PNG/SVG/PDF \\
Bead colourmap              & coolwarm & Any of 30+ matplotlib colourmaps \\
Graph accent colour         & Royal Blue & 24 named colour choices \\
Show graph titles           & on     & Toggle titles above plots \\
Show grid                   & on     & Toggle gridlines \\
Show bead labels            & on     & Residue letters on bead models ($\leq$60\,aa) \\
Transparent background      & on     & Transparent PNG/SVG export \\
UI font size                & 12\,pt & Global font size (8--24\,pt) \\
Dark theme                  & off    & Light/dark toggle \\
Enable tooltips             & on     & Hover tooltips on Settings widgets \\
\end{longtable}

% ============================================================
\section{Use cases}
\label{sec:usecases}

\subsection{Case 1: rapid characterisation of a newly cloned IDP}

A researcher clones a 320-residue disordered protein.
They paste the FASTA sequence into BEER; within seconds the Analysis section
reports FCR\,=\,0.24, NCPR\,=\,$-$0.08 (polyampholyte), $\kappa$\,=\,0.31
(moderate charge segregation), disorder fraction\,=\,0.62, and
LLPS score\,=\,0.71 (high propensity).
The Graphs section shows a flat hydrophobicity profile consistent with
disorder, a charge-decoration Das--Pappu plot placing the sequence in the
weak polyampholyte quadrant, and a saturation mutagenesis heatmap identifying
a cluster of hydrophobic residues at positions 140--165 as the principal
drivers of GRAVY.
The researcher fetches the AlphaFold structure and inspects the pLDDT profile
confirming low confidence ($<$50) throughout the central disordered region.
The ELM fetch returns three validated SLiM instances, one of which overlaps
a PKA phosphorylation site predicted by the built-in scanner.
The complete analysis is exported as a PDF report.

\subsection{Case 2: comparison of two homologous membrane proteins}

Starting from PDB~1GP2 (bovine opsin) the researcher imports the PDB file,
triggering automatic extraction of all protein chains and auto-analysis of
the first chain.
The TM topology plot shows seven predicted transmembrane helices with
inside-positive topology consistent with the known GPCR architecture.
The researcher then opens the Compare tab, enters a second sequence (a closely
related class-A GPCR), and obtains a side-by-side table of all physicochemical
properties, immediately identifying differences in pI, aromatic fraction, and
LLPS score between the two proteins.

\subsection{Case 3: aggregation risk assessment of a therapeutic antibody}

A biopharmaceutical scientist analyses the variable heavy chain of a candidate
monoclonal antibody.
The aggregation module identifies two ZYGGREGATOR hotspots in the CDR regions
(positions 31--36 and 97--104) with PASTA energies of $-0.82$ and $-0.94$,
respectively, flagging these as high-risk for aggregation.
The CamSol mean score of $-0.07$ (borderline) and the CamSol profile plot
highlight the same CDR stretches as low-solubility windows.
These quantitative outputs directly inform a mutagenesis campaign to improve
antibody solubility.

% ============================================================
\section{Comparison with existing tools}
\label{sec:comparison}

Table~\ref{tab:comparison} contrasts BEER with the most widely used
comparable tools.

\begin{longtable}{p{2.8cm}p{1.8cm}p{1.8cm}p{1.8cm}p{1.8cm}p{1.8cm}}
\caption{Feature comparison of BEER with selected related software.}
\label{tab:comparison}\\
\toprule
\textbf{Feature} & \textbf{BEER} & \textbf{Prot-Param} & \textbf{IUPRED3} & \textbf{CAMELOT} & \textbf{localCIDER} \\
\midrule
\endfirsthead
\toprule
\textbf{Feature} & \textbf{BEER} & \textbf{Prot-Param} & \textbf{IUPRED3} & \textbf{CAMELOT} & \textbf{localCIDER} \\
\midrule
\endhead
\bottomrule
\endlastfoot
Basic properties     & \checkmark & \checkmark & \checkmark & --- & --- \\
Disorder profile     & \checkmark & ---        & \checkmark & \checkmark & \checkmark \\
Charge analysis      & \checkmark & partial    & ---        & \checkmark & \checkmark \\
LLPS score           & \checkmark & ---        & ---        & partial    & partial \\
Aggregation          & \checkmark & ---        & ---        & ---        & --- \\
Solubility           & \checkmark & ---        & ---        & ---        & --- \\
PTM prediction       & \checkmark & ---        & ---        & ---        & --- \\
Signal peptide       & \checkmark & ---        & ---        & ---        & --- \\
Amphipathic helix    & \checkmark & ---        & ---        & ---        & --- \\
SCD                  & \checkmark & ---        & ---        & ---        & \checkmark \\
RNA-binding          & \checkmark & ---        & ---        & ---        & --- \\
Tandem repeats       & \checkmark & ---        & ---        & ---        & --- \\
BLAST                & \checkmark & ---        & ---        & ---        & --- \\
Structure viewer     & \checkmark & ---        & ---        & ---        & --- \\
AlphaFold fetch      & \checkmark & ---        & ---        & ---        & --- \\
Pfam domains         & \checkmark & ---        & ---        & ---        & --- \\
ELM database         & \checkmark & ---        & ---        & ---        & --- \\
DisProt lookup       & \checkmark & ---        & ---        & ---        & --- \\
PhaSepDB lookup      & \checkmark & ---        & ---        & ---        & --- \\
Multi-sequence/chain & \checkmark & ---        & ---        & ---        & --- \\
Offline operation    & \checkmark & ---        & ---        & \checkmark & \checkmark \\
Interactive plots    & 31         & 0          & 1          & few        & few \\
GUI                  & desktop    & web        & web        & web        & Python API \\
Open source          & GPLv2      & ---        & web-only   & ---        & MIT \\
\end{longtable}

% ============================================================
\section{Validation}

Core computational functions are validated by unit tests in
\texttt{tests/test\_beer.py} covering: composition counts, net charge at
multiple pH values, FCR/NCPR/kappa calculations, Shannon entropy, GRAVY,
TM helix prediction (known single-pass and polytopic membrane proteins),
LARKS detection, LLPS scoring, PTM regex scanning, signal peptide cleavage
sites, amphipathic helix moment, SCD, RBP score, and tandem repeat
identification.
All tests pass on Python 3.11 and 3.12 on Linux (CentOS~Stream~8), macOS
(Sequoia~15.3), and Windows 11.

For qualitative validation, BEER was applied to three well-characterised
proteins: (i) human p53 (UniProt P04637) to verify disorder profile, PTM
sites (phosphorylation, acetylation, ubiquitination), and LLPS score
consistent with its known ability to form condensates; (ii) GCN4 leucine
zipper peptide (PDB~2ZTA) to confirm coiled-coil prediction and correct
amphipathic moment; (iii) bovine rhodopsin (PDB~1GP2) to verify the seven-TM
topology prediction with correct inside-positive orientation.

% ============================================================
\section{Availability and requirements}
\label{sec:avail}

\begin{description}[noitemsep]
  \item[Repository] \url{https://github.com/chemgame/BEER}
  \item[Licence] GNU General Public Licence v2.0
  \item[Language] Python 3.11--3.12
  \item[Platforms] Windows 10/11, macOS 12+, Linux (tested on CentOS~Stream~8)
  \item[Dependencies] BioPython $\geq$1.80, NumPy $\geq$1.24,
        Matplotlib $\geq$3.7, PyQt5 $\geq$5.15, mplcursors;
        PyQtWebEngine (optional, for 3-D viewer)
  \item[Installation] \texttt{pip install biopython matplotlib pyqt5 mplcursors}
  \item[Launch] \texttt{python beer.py}
\end{description}

% ============================================================
\section{Discussion and conclusions}

BEER addresses a clear gap in the bioinformatics toolkit: the need for a
desktop application that combines, in a single self-contained and open-source
program, the breadth of analytical coverage previously requiring visits to a
dozen or more web servers.
Its design philosophy---a single executable Python file, no proprietary
dependencies, local computation for all core analyses---makes it suitable
for use in secure environments, on air-gapped workstations, and in teaching
contexts where students need immediate, reproducible results without managing
API credentials.

Several of the implemented methods deserve specific comment.
The ZYGGREGATOR/CamSol/PASTA aggregation triad provides three independently
derived and complementary views of aggregation risk: the Zyggregator propensity
reflects the intrinsic $\beta$-aggregation tendency of the primary sequence,
the PASTA diagonal energies measure amyloidogenic self-pairing strength, and
CamSol quantifies the net balance between solubilising and aggregation-prone
physicochemical features.
Used together they provide actionable guidance for protein engineering without
the need for time-consuming experimental assays.

The integration of SCD and charge-distribution analysis is particularly
relevant to the rapidly growing IDP and condensate biology fields.
SCD and $\kappa$ together characterise both the magnitude and the spatial
patterning of sequence charges, features that are increasingly recognised as
the primary determinants of IDP conformational behaviour and condensate
material properties~\citep{pappu2023}.

Future development directions include: (1) integration of sequence-based LLPS
predictor machine-learning models; (2) a scripting API for batch analysis of
proteomes; (3) native packaging as a standalone executable for each supported
platform; and (4) extension of the structural analysis module to support
molecular dynamics trajectory files.

% ============================================================
\section*{Acknowledgements}

The author thanks Gerhard Hummer for scientific discussions and support at
the Max Planck Institute of Biophysics.
Development of BEER was assisted in part by large language model tools.

% ============================================================
\section*{Funding}

S.M.\ acknowledges support from the Max Planck Society.

% ============================================================
\section*{Conflict of interest}

None declared.

% ============================================================
\begin{thebibliography}{99}

\bibitem[Almagro~Armenteros et~al.(2019)]{almagro2019}
Almagro~Armenteros, J.J., Tsirigos, K.D., S{\o}nderby, C.K., Petersen, T.N.,
Winther, O., Brunak, S., von~Heijne, G., Nielsen, H. (2019).
SignalP~5.0 improves signal peptide predictions using deep neural networks.
\textit{Nature Biotechnology}, \textbf{37}, 420--423.

\bibitem[Chou and Fasman(1974)]{chou1974}
Chou, P.Y. and Fasman, G.D. (1974).
Prediction of protein conformation.
\textit{Biochemistry}, \textbf{13}, 222--245.

\bibitem[Cock et~al.(2009)]{cock2009}
Cock, P.J.A., Antao, T., Chang, J.T., Chapman, B.A., Cox, C.J., Dalke, A.,
Friedberg, I., Hamelryck, T., Kauff, F., Wilczynski, B. and de~Hoon, M.J.L.
(2009).
Biopython: freely available Python tools for computational molecular biology
and bioinformatics.
\textit{Bioinformatics}, \textbf{25}, 1422--1423.

\bibitem[Das and Pappu(2013)]{das2013}
Das, R.K. and Pappu, R.V. (2013).
Conformations of intrinsically disordered proteins are influenced by linear
sequence distributions of oppositely charged residues.
\textit{Proceedings of the National Academy of Sciences}, \textbf{110},
13392--13397.

\bibitem[Das et~al.(2015)]{das2015}
Das, R.K., Ruff, K.M. and Pappu, R.V. (2015).
Relating sequence encoded information to form and function of intrinsically
disordered proteins.
\textit{Current Opinion in Structural Biology}, \textbf{32}, 102--112.

\bibitem[Dosztanyi et~al.(2005)]{dosztanyi2005}
Dosztanyi, Z., Csizmok, V., Tompa, P. and Simon, I. (2005).
IUPred: web server for the prediction of intrinsically unstructured regions
of proteins based on estimated energy content.
\textit{Bioinformatics}, \textbf{21}, 3433--3434.

\bibitem[Eisenberg et~al.(1984)]{eisenberg1984}
Eisenberg, D., Schwarz, E., Komaromy, M. and Wall, R. (1984).
Analysis of membrane and surface protein sequences with the hydrophobic
moment plot.
\textit{Journal of Molecular Biology}, \textbf{179}, 125--142.

\bibitem[Gasteiger et~al.(2005)]{gasteiger2005}
Gasteiger, E., Hoogland, C., Gattiker, A., Duvaud, S., Wilkins, M.R.,
Appel, R.D. and Bairoch, A. (2005).
Protein identification and analysis tools on the ExPASy server.
In: Walker, J.M. (ed.), \textit{The Proteomics Protocols Handbook},
Humana Press, pp.\ 571--607.

\bibitem[Hughes et~al.(2018)]{hughes2018}
Hughes, M.P., Sawaya, M.R., Boyer, D.R., Goldschmidt, L., Rodriguez, J.A.,
Cascio, D., Chong, L., Gonen, T. and Eisenberg, D.S. (2018).
Atomic structures of low-complexity protein segments reveal kinked
$\beta$ sheets that assemble networks.
\textit{Science}, \textbf{359}, 698--701.

\bibitem[Kumar et~al.(2020)]{kumar2020}
Kumar, M., Gouw, M., Michael, S., Samano-Sanchez, H., Pancsa, R.,
Glavina, J., Diakogianni, A., Bhowmick, A., Bhatt, D., Naaktgeboren, M.,
Meszaros, B., Davey, N.E., Dosztanyi, Z., Gibson, T.J. and Dinkel, H.
(2020).
ELM---the eukaryotic linear motif resource in 2020.
\textit{Nucleic Acids Research}, \textbf{48}, D296--D306.

\bibitem[Kyte and Doolittle(1982)]{kyte1982}
Kyte, J. and Doolittle, R.F. (1982).
A simple method for displaying the hydropathic character of a protein.
\textit{Journal of Molecular Biology}, \textbf{157}, 105--132.

\bibitem[Lancaster and Bhatt(2014)]{lancaster2014}
Lancaster, A.K. and Bhatt, D. (2014).
Computational prediction of prion-like domains.
\textit{Methods in Molecular Biology}, \textbf{1216}, 113--128.

\bibitem[Mukherjee(2025)]{mukherjee2025}
Mukherjee, S. (2025).
BEER: Biochemical Estimator \& Explorer of Residues --- a comprehensive
software suite for protein sequence analysis.
\textit{arXiv}:2504.20561.
DOI: \href{https://doi.org/10.48550/arXiv.2504.20561}{10.48550/arXiv.2504.20561}.

\bibitem[Pappu et~al.(2023)]{pappu2023}
Pappu, R.V., Cohen, S.R., Dar, F., Farag, M. and Kar, M. (2023).
Phase transitions of associative biomacromolecules.
\textit{Chemical Reviews}, \textbf{123}, 8945--8987.

\bibitem[Romero et~al.(2001)]{romero2001}
Romero, P., Obradovic, Z., Li, X., Garner, E.C., Brown, C.J. and
Dunker, A.K. (2001).
Sequence complexity of disordered protein.
\textit{Proteins: Structure, Function, and Bioinformatics}, \textbf{42},
38--48.

\bibitem[Sawle and Ghosh(2015)]{sawle2015}
Sawle, L. and Ghosh, K. (2015).
A theoretical method to compute sequence dependent configurational properties
in charged polymers and proteins.
\textit{The Journal of Chemical Physics}, \textbf{143}, 085101.

\bibitem[Sormanni et~al.(2015)]{sormanni2015}
Sormanni, P., Aprile, F.A. and Vendruscolo, M. (2015).
The CamSol method of rational design of protein mutants with enhanced
solubility.
\textit{Journal of Molecular Biology}, \textbf{427}, 478--490.

\bibitem[Tartaglia and Vendruscolo(2008)]{tartaglia2008}
Tartaglia, G.G. and Vendruscolo, M. (2008).
The Zyggregator method for predicting protein aggregation propensities.
\textit{Chemical Biology}, \textbf{15}, 1008--1018.

\bibitem[Trovato et~al.(2007)]{trovato2007}
Trovato, A., Chiti, F., Maritan, A. and Seno, F. (2007).
Insight into the structure of amyloid fibrils from the analysis of globular
proteins.
\textit{PLoS Computational Biology}, \textbf{3}, e17.

\end{thebibliography}

\end{document}
